\chapter*{Conferencia X}
\addcontentsline{toc}{chapter}{LECTURA I - El timbre y el carácter del violín}
Señores: Hay violinistas que prefieren esta desigualdad en la potencia del sonido. En mi experiencia, una excesiva prominencia del sonido de la cuerda de sol no resulta del todo agradable al oído, especialmente en el violín solista. Considero este punto lo suficientemente importante como para reiterar lo dicho anteriormente: la hermosa armonía que se logra con cuerdas equilibradas, tocadas en acordes dobles, es el principal atractivo de la música para violín solista.
(6) . El poste es a la vez lo más irritante, poderoso, indispensable y aparentemente inocente dentro o fuera del violín, alrededor o cerca de él. No se cae tan a menudo como nos hace caer. Si tuviera la capacidad de reír, sería el principal "mono". Para el violinista, el poste es el principal objeto de preocupación. Su poder para el bien es angelical. Su poder para el mal es satánico. ¿Puede este objeto aparentemente inocente tener el poder de disminuir tanto el volumen como la intensidad del sonido del violín?
¡Que caiga ahí abajo! Porque yo
Disfruto de frecuentes tomas de este ángel-diablo, por lo tanto, como tema, no se agotará de inmediato. En este momento mencionaré solo tres puntos relacionados con la publicación que sirven para disminuir el volumen y la intensidad del tono:
(a) . Masa del poste. (b) . Posición del poste. (c) . Longitud del poste. En la condición normal de la salida derecha, el mayor efecto sobre el volumen y la intensidad del tono,
Debido a la gran cantidad de correo, no se puede demostrar. 160	En-

La ampliación de esta salida permite la entrada de un poste de gran masa. Si bien la masa del poste no disminuye la potencia del sonido en la misma medida que otros factores, dicha disminución es perceptible. No conozco ninguna regla fija que determine la masa del poste. Al igual que muchos otros factores que afectan el sonido del violín, la masa óptima para cada violín solo puede determinarse mediante prueba y error.
(b). La posición del poste ejerce mayor influencia sobre la potencia tonal que la masa del poste; es decir, sobre cualquier masa que haya probado. Recuerdo haber leído en mi juventud que la posición correcta del poste es una pulgada por debajo del pie derecho del puente. Más adelante, aprendí a desaprender esa lección. Tuve que aprender que la posición del poste varía según las variaciones en la rigidez de la caja de resonancia. Todos los efectos
sobre el tono, debido al poste, se manifiestan principalmente en las cuerdas A y E; y tales efectos pueden dirigirse en gran medida a cualquiera de estas cuerdas por la posición del poste. Cuando se encuentra esa posición que soporta por igual las cuerdas A y E, entonces, mover el poste hacia la izquierda opera para disminuir la potencia tonal de E; y, por el contrario, mover el poste hacia la derecha, disminuye la potencia tonal de A. Nuevamente, mover el poste hacia abajo, opera para disminuir la potencia tonal tanto de A como de E. Nuevamente, colocar el poste sobre el puente opera para disminuir la potencia tonal de las cuerdas A y E en una cantidad igual al refuerzo debido a la acción simpática entre estas cuerdas y esa parte de la caja de resonancia debajo de ellas. Colocar el
161

El poste situado encima del puente transfiere la acción de la caja de resonancia a la mitad inferior de la misma; por el contrario, colocar el poste debajo del puente produce la acción de la caja de resonancia en la mitad superior. Deseo que se entienda que me refiero únicamente a la acción de la caja de resonancia provocada por la acción de las cuerdas A y E; además, que en cualquiera de estas dos posiciones del poste no se produce ningún cambio perceptible en el tono de las cuerdas G y E.
Cuerdas D.
[En una ocasión posterior aparecerá una demostración práctica del hecho de que el cuadrante inferior derecho de la caja de resonancia no produce ondas sonoras cuando el poste se coloca debajo del pie derecho del puente.]
Colocar el poste directamente debajo del pie derecho del puente reduce la potencia del sonido. En esta posición, es evidente que el golpe de las cuerdas se distribuye por igual entre ambas placas; sin embargo, en mi opinión, la caja de resonancia sigue recibiendo el mayor impacto en el aire contenido debido a su mayor proximidad a las cuerdas. Esta interpretación se basa en el hecho de que la acción simpática disminuye con el cuadrado de la distancia, o dicho de otro modo, la proximidad aumenta la acción simpática.
[Para el regulador de tono del violín, existen experiencias con la humanidad que poseen un interés más que pasajero. De hecho, algunas de esas experiencias dejan una impresión en la memoria tan imborrable como la impresión en la "madera quemada". Por lo tanto: Después de haber ajustado cuidadosamente la caja de resonancia
162

espesores para responder a la acción de las cuerdas de calibre 2, después de que las superficies interiores se hayan preparado cuidadosamente, después de que el área de salidas haya recibido atención, después de que el diapasón se haya preparado y ajustado cuidadosamente, después de seleccionar y probar diferentes densidades en diferentes puentes, después de seleccionar las cuerdas más selectas, después de determinar la posición exacta para el puente y el poste, después de probar el tono para uniformidad de potencia en dos octavas en cada cuerda, probando la calidad de los tonos de doble cuerda, probando el brillo del tono, probando los tonos de pizzicato, probando armónicos simples, armónicos melódicos, armónicos a bassa, después de declarar que su trabajo se ha completado al límite de su capacidad:
Entra el señor Addlepate.
(Sr. A.) "Estoy buscando un violín de primera clase para uso personal. No lo sabía, pero podría encontrar uno en su establecimiento."
Con bastante seguridad, le entregas un violín terminado, diciéndole que este es prácticamente lo mejor que puedes hacer.
(Sr. A.)	¿Puedo llevármelo a casa y probarlo?
Cuídalo bien."	
"Ciertamente."	
Después de uno o dos meses, vuelva a entrar el Sr. A.	
(Sr. A.)	"Diga, este violín no tiene solo el
tono que quiero	Creo que Sam Jones tiene uno.
un poco mejor que el tuyo."
Al mirar ese violín, tu respiración se convierte en un jadeo. Tu alma cuidadosamente ajustada ha desaparecido. De pie hacia el centro, e inclinado hacia atrás, hay un poste de roble blanco, tallado con una navaja desafilada; una muesca profunda está cortada alrededor de...
163

el extremo superior; en esa muesca hay una cuerda de algodón sucia doblada y retorcida, atada firmemente con un nudo cuadrado; cada extremo de la cuerda, que sale por cualquiera de las aberturas, se lleva por detrás de la espalda, donde se ata con un nudo de lazo; de ahí cuelga 1/4 de yarda de guirnalda. "¿Qué te hace jadear tanto?" "¿No puede ser ese poste de roble blanco?"
"¡No!"
"¿Posiblemente sea esa guirnalda?"
Ante semejante provocación, solo hay una manera de manifestar patriotismo. Esa manera consiste en propinar una patada entusiasta directamente al cerebro del señor Addlepate. Es imposible no darse cuenta.
c) La longitud del mástil, cuando es tan grande que dobla la caja de resonancia hacia arriba, disminuye el volumen y la duración de las cuerdas A y E, pero no su intensidad en la misma medida. De hecho, un mástil demasiado largo suele aumentar la intensidad del sonido de estas cuerdas. Algunos violinistas, especialmente los primeros violinistas en posiciones donde la propagación de las ondas sonoras es difícil, consideran valiosa esta mayor intensidad del sonido de las cuerdas A y E. Un mástil demasiado largo también disminuye la potencia del sonido de las cuerdas G y D.
De esta manera, el puesto manifiesta su poder para el bien y
para el mal.
(7. ) El puente es un factor importante para disminuir tanto el volumen como la intensidad del sonido; sin embargo, su mayor influencia se manifiesta en el volumen. El puente posee siete características de gran interés para el estudiante de violín:
164

(a). (b). (c). (d). (e).	Altura. Grosor. Distancia entre vástagos o pedestales. Densidad de la fibra. Trabajo de volutas.
(f). Madurez de la madera.	
(g) .	Posiciónate en la mesa de resonancia.
(a) . La altura del puente puede ser tan alta o tan baja que disminuya tanto el volumen como la intensidad del tono. No puedo resolver el problema del puente sin una solución parcial para el problema del diapasón. Para mayor claridad, primero es necesario considerar el diapasón en lo que respecta a la proximidad a la caja de resonancia. Si el diapasón no se extendiera sobre la caja de resonancia, entonces el problema de la altura del puente podría resolverse por sí solo. Esa parte del diapasón que se extiende sobre la caja de resonancia puede disminuir tanto el volumen como la intensidad del tono. Este hecho se puede demostrar fácilmente. Así: En un violín, con una altura de arco igual a f pulgadas, bajar el diapasón a 1 pulgada de la caja de resonancia produce un tono débil y delgado, incluso con una disminución correspondiente de la altura del puente. Como veo este fenómeno, la razón de tal pérdida de potencia tonal es la siguiente: Bajar el diapasón hace que su línea de superficie se acerque a una paralela a la línea de la caja de resonancia. Es evidente que, si estas dos líneas fueran perfectamente paralelas, todas las ondas sonoras que se originan debajo del diapasón se reflejarían directamente sobre la caja de resonancia; y, debido a que solo viajan
165

A corta distancia, las ondas sonoras reflejadas deben impactar la caja de resonancia con su fuerza inicial; una fuerza lo suficientemente grande como para disminuir la amplitud de la oscilación de la caja de resonancia. Es evidente que esta disminución de amplitud reduce la fuerza de los golpes subsiguientes que la caja de resonancia produce sobre el aire contenido. De ahí la pérdida de potencia tonal.
El problema del diapasón radica en la línea de su superficie inferior y en la colocación de dicha línea a la menor distancia posible de la caja de resonancia, evitando al mismo tiempo que las ondas sonoras reflejadas incidan sobre la caja de resonancia de tal manera que disminuyan la amplitud de su oscilación.
Como ya se mencionó, considero que las mejores líneas en la superficie inferior del diapasón son una línea recta en el fondo del hueco, y que este se extiende bien hacia la base del mástil. Este último punto debe estar determinado por la graduación de la tapa armónica y su curvatura. Así, cuando la graduación sitúa la parte más delgada de la tapa armónica a medio camino entre la posición del puente y el extremo superior de la placa, el hueco del diapasón no debe extenderse hasta la base del mástil, sino hasta un punto situado a 5 cm (2 pulgadas), o un poco más, de la base.
[En una fecha posterior, esta forma para la superficie inferior del diapasón, junto con esta forma para la graduación de la caja de resonancia, ligeramente modificada, se describirá en la producción de máximo
potencia tonal.]
Es claramente evidente que la posición del cuello
166

Puede operar para colocar el diapasón cerca o lejos de la caja de resonancia. También es evidente que cuando el punto de máxima oscilación de la caja de resonancia está más cerca del extremo de la placa, mayor debe ser la distancia a la superficie inferior del diapasón, porque cuanto más larga sea la cavidad, menor será la oblicuidad de su línea y más directo será el retorno de las ondas sonoras a la caja de resonancia.
Aún queda por considerar las placas arqueadas. En estos casos, la línea de la superficie inferior del diapasón y su distancia a la caja de resonancia deben ser justo las opuestas; el extremo inferior del diapasón debe acercarse mucho más a la caja de resonancia, mientras que la superficie inferior debe ser plana. A primera vista, estos hechos parecen increíbles. Confieso que me asombré al ver y tocar por primera vez un violín Carl Johann Flicker. Como saben, este violín está construido sobre unas cinturas enormes. El hueco en el extremo inferior del diapasón (calculando) estaba a menos de una pulgada de la caja de resonancia. La altura del puente se redujo en consecuencia. Antes de aplicar el arco, esperaba oír un tono débil y fino. Pero su volumen era bastante grande y la intensidad parecía ser la media para el arco. No tuve la oportunidad de determinar la altura del arco de las placas del Flicker. Pude ver que el arco largo estaba distribuido de manera bastante uniforme desde los extremos de las placas hasta la posición del puente. Al estudiar esta situación, se hace evidente que para elevar el extremo inferior de este

Elevar el diapasón a una altura que permita que su superficie inferior refleje las ondas sonoras hacia el puente es poco práctico, no imposible, simplemente poco práctico; y porque la prodigiosa altura de las cuerdas dificultaría demasiado la ejecución. El Flicker no era difícil de tocar.
Así se demuestra que la altura del puente debe depender de la altura del diapasón, y la altura del diapasón debe depender de la altura del arco. En el caso del violín Flicker, la superficie inferior plana del diapasón, que descansa sobre un mástil alto, dirigía las ondas sonoras hacia el mástil, y aparentemente, este es el único diseño práctico para violines con la mayor altura de arco. Las salidas del Flicker eran inusualmente grandes y parecían estar inusualmente juntas. El resorte del arco comenzaba directamente dentro del fileteado. A estos hechos atribuyo la inusual potencia tonal de este violín.
violín en comparación con la potencia tonal de otros violines de su clase.
Es evidente que aumentar el ancho del diapasón incrementa el número de ondas sonoras reflejadas hacia la caja de resonancia; por lo tanto, dicho aumento de ancho disminuye la potencia tonal. El diapasón es necesario, sin embargo, incluso con el mejor ajuste posible, el diapasón disminuye la potencia tonal del violín. Así pues, el ancho, la línea reflectante inferior y la distancia de dicha línea a la caja de resonancia son potencialidades del diapasón que requieren la mayor precisión posible.
Atención del regulador de tono del violín. 168	Falta

El descuido conlleva sanciones. En lo que respecta a la distancia entre la caja de resonancia y la línea de la superficie inferior del diapasón, no se pueden aplicar reglas estrictas, ya que la curvatura y la graduación de la caja de resonancia rigen la norma.
La altura deseada de las cuerdas desde el diapasón y la posición del puente sobre la caja de resonancia determinan la altura del puente. Se puede permitir cierta flexibilidad en la altura de las cuerdas sobre el diapasón para adaptarse a diferentes gustos. Mi elección personal para esta altura es clara en todo momento. Algunos prefieren una mayor altura para el Sol y una menor altura para el Mi, mientras que otros prefieren una mayor altura en todas las cuerdas. Una vez conocí a un violinista experto, con una longitud inusual de manos y dedos, que colocaba las cuerdas una pulgada por encima del extremo inferior del diapasón. Le pregunté por qué. Él respondió: "Por dos razones;
Una de las razones es que mi violín produce un sonido más potente; la otra, que los tonos pizzicato son más nítidos. No creo que su primera razón sea válida en todos los casos; al contrario, he observado que al tensar las cuerdas de esta manera se debilita la potencia del sonido.
Un tono pizzicato claro es muy deseable. El tono pizzicato entrecortado debido a la oscilación de las cuerdas que recibe interferencia del diapasón es algo inadmisible. Para asegurar un tono pizzicato claro, manteniendo las cuerdas a una distancia de , o menos, del extremo inferior del diapasón, recurro al siguiente tratamiento de la parte superior del diapasón.
169

superficie. Con un pequeño cepillo de carpintero, voy rebajando esta superficie hasta que su línea presente una ligera curva de un extremo a otro. Se retira la menor cantidad de madera debajo de la cuerda E, y la mayor cantidad debajo de la cuerda G, debido a que la amplitud de la oscilación de la cuerda aumenta de E a G. El punto más bajo de esta curva, determinado por la regla, no se sitúa más arriba de C en alt, debido a que acortar una cuerda reduce la amplitud de su oscilación. El análisis del puente se retomará en la próxima hora.
